\section{Count Subsequences with Sum K}
\label{sec:rec-count-subsequence-sum-k}

\problemstatement{Given an array $\textit{arr}$ and an integer $k$, count
\textbf{how many} subsequences of $\textit{arr}$ sum to exactly $k$.}

\begin{patternbox}
The third variation on Section~\ref{sec:rec-power-set}'s pick/not-pick
tree: from \emph{``generate all''} (Section~\ref{sec:rec-power-set}) to
\emph{``does one exist''} (Section~\ref{sec:rec-check-subsequence-sum-k})
to, now, \emph{``count how many.''} Counting needs neither full
generation nor early short-circuiting -- it needs both branches'
\textbf{counts} added together.
\end{patternbox}

\subsection*{Understanding the problem carefully}

Every complete pick/not-pick decision sequence produces exactly one
subsequence, and the base case can directly check whether \emph{that}
subsequence's sum equals $k$, contributing $1$ to the count if so, $0$
otherwise. Since the pick and not-pick branches partition all
subsequences into two disjoint groups (those that include
$\textit{arr}[\textit{index}]$ and those that don't), the total count is
simply the \textbf{sum} of the counts returned by both branches -- no
branch can be skipped, since both might contain qualifying subsequences.

\begin{ideabox}[Sum the Counts from Both Branches]
$\textsc{CountSum}(\textit{index}, \textit{sum})$: if
$\textit{index} = n$: return $1$ if $\textit{sum} = k$, else $0$.
Otherwise: return $\textsc{CountSum}(\textit{index}+1, \textit{sum} +
\textit{arr}[\textit{index}]) + \textsc{CountSum}(\textit{index}+1,
\textit{sum})$.
\end{ideabox}

\subsection*{Why summing (not short-circuiting) is required here}

Unlike Section~\ref{sec:rec-check-subsequence-sum-k}, where finding one
success already fully answers the yes/no question, here a successful
pick branch tells us nothing about how many \emph{additional} successes
might be hiding in the not-pick branch -- both must be fully explored and
their individual counts combined, since the two branches' qualifying
subsequences are entirely different (disjoint) subsequences, both of
which need to be counted toward the final total.

\subsection*{Algorithm}

\begin{enumerate}
  \item $\textsc{CountSum}(\textit{index}, \textit{sum})$:
  \begin{enumerate}
    \item If $\textit{index} = n$: return $1$ if $\textit{sum}=k$, else
          $0$.
    \item Return $\textsc{CountSum}(\textit{index}+1, \textit{sum} +
          \textit{arr}[\textit{index}]) +
          \textsc{CountSum}(\textit{index}+1, \textit{sum})$.
  \end{enumerate}
  \item Call $\textsc{CountSum}(0,0)$.
\end{enumerate}

\subsection*{Dry run}

\dryrun{Take $\textit{arr} = [1,2,1]$, $k=2$.}

\begin{itemize}
  \item $\textsc{CountSum}(0,0)$: pick $1$:
        $\textsc{CountSum}(1,1)$; not-pick: $\textsc{CountSum}(1,0)$.
  \item $\textsc{CountSum}(1,1)$ (came from picking index $0$): pick
        $2$: $\textsc{CountSum}(2,3)$; not-pick:
        $\textsc{CountSum}(2,1)$.
  \begin{itemize}
    \item $\textsc{CountSum}(2,3)$: pick $1$: base, sum$=4\ne2$,
          return $0$. Not-pick: base, sum$=3\ne2$, return $0$. Total
          $0$.
    \item $\textsc{CountSum}(2,1)$: pick $1$: base, sum$=2=2$! return
          $1$. Not-pick: base, sum$=1\ne2$, return $0$. Total $1$.
  \end{itemize}
  \item $\textsc{CountSum}(1,1)=0+1=1$.
  \item $\textsc{CountSum}(1,0)$ (came from not-picking index $0$):
        pick $2$: $\textsc{CountSum}(2,2)$; not-pick:
        $\textsc{CountSum}(2,0)$.
  \begin{itemize}
    \item $\textsc{CountSum}(2,2)$: pick $1$: base, sum$=3\ne2$,
          return $0$. Not-pick: base, sum$=2=2$! return $1$. Total $1$.
    \item $\textsc{CountSum}(2,0)$: pick $1$: base, sum$=1\ne2$,
          return $0$. Not-pick: base, sum$=0\ne2$, return $0$. Total
          $0$.
  \end{itemize}
  \item $\textsc{CountSum}(1,0)=1+0=1$.
  \item $\textsc{CountSum}(0,0)=1+1=2$.
\end{itemize}

The answer is $2$: the qualifying subsequences are $\{1,1\}$ (indices
$0,2$) and $\{2\}$ (index $1$ alone).

\subsection*{C++ implementation}

\begin{lstlisting}
#include <bits/stdc++.h>
using namespace std;

int countSum(int index, int sum, vector<int>& arr, int k) {
    if (index == (int)arr.size()) return sum == k ? 1 : 0;

    int pick = countSum(index + 1, sum + arr[index], arr, k);
    int notPick = countSum(index + 1, sum, arr, k);
    return pick + notPick;
}

int main() {
    vector<int> arr = {1, 2, 1};
    cout << "Count: " << countSum(0, 0, arr, 2) << endl;
    return 0;
}
\end{lstlisting}

\begin{complexitybox}
\textbf{Time Complexity.} $O(2^n)$ -- every one of the $2^n$
subsequences must be considered, since any of them might contribute to
the count.
\textbf{Space Complexity.} $O(n)$ for the recursion stack.
\end{complexitybox}

\begin{takeawaybox}
The three variations seen so far on the same pick/not-pick tree --
generate everything (Section~\ref{sec:rec-power-set}), short-circuit on
the first success (Section~\ref{sec:rec-check-subsequence-sum-k}), and
sum counts from both branches (this section) -- are worth holding side by
side as a single mental checklist: \emph{generate, exists, or count} is
usually the very first distinction to make about any new backtracking
problem, since it determines how the two branches' results should be
combined at every level.
\end{takeawaybox}
